\documentclass[11pt,a4paper]{letter}
\usepackage[T1]{fontenc}
\usepackage[utf8]{inputenc}
\usepackage{lmodern}
\usepackage[margin=1.1in]{geometry}
\usepackage[hidelinks]{hyperref}

\signature{Tommaso R. Marena}
\address{Tommaso R. Marena\\ \texttt{marenatommaso@gmail.com}}
\begin{document}

\begin{letter}{The Editors\\ \emph{Nature Methods}}

\opening{Dear Editors,}

I am submitting ``A benchmark cannot order more methods than its labels can
distinguish'' for consideration as an Article.

Community benchmarks are the field's instrument for deciding which methods
work, and the ordering they publish is treated as data. This manuscript proves
that it usually is not. A benchmark whose annotations are wrong at rate
$\varepsilon$ can place at most $\lceil 1/2\varepsilon \rceil$ methods in a
certifiable order --- however many enter, and however the data are analysed ---
and past that count the ordering is \emph{not a function of the observations}:
for two close methods we construct two ground truths, each consistent with
everything the benchmark recorded, one favouring each. This is an impossibility
result, not a power calculation, and no amount of resampling touches it.

Measuring the annotation error rate directly, from repeat determinations of the
same protein in MobiDB, gives the Critical Assessment of Intrinsic Disorder a
capacity of seven. It admits 117 entrants.

The paper does not stop at the critique. The bound constrains \emph{labels},
while a ranking statistic is a function of ordered \emph{pairs}, and a pair is
reversed only when both its residues are mis-annotated in opposite directions
--- an exactly second-order event. We give the identity, measure both rates,
specify a pairwise protocol that raises the capacity from 8 to 51, re-score the
whole CAID3 field under it, and show the resulting ranking reproduces itself on
a held-out round as well as the published one does. CAID4's references are not
yet public, so this is adoptable rather than retrospective.

Three features may interest \emph{Nature Methods} specifically. First, every
theorem is machine-checked in Lean~4, \texttt{sorry}-free and on standard
axioms; formalising the central identity refuted two numbers we had been
publishing, including our own headline capacity, and both corrections are
documented in the manuscript rather than quietly absorbed. Second, we give
disorder prediction what we believe is its first distribution-free operating
guarantee: fix a tolerated miss rate and conformal risk control returns a
threshold that achieves it on an unseen protein, at a price that ranges from
32\% to 97\% of the protein flagged across the CAID3 field --- a difference AUC
cannot see. Third, the same design argument governs reporting: because the
dependence correction a screen must pay is set by how many hypotheses it states,
testing regions rather than residues yields 14.5$\times$ the discoveries at the
same guaranteed false discovery rate, which we prove and then demonstrate.

The accompanying predictor, DisorderNet, ranks first of 115 entrants on three of
five CAID3 references and first among full-coverage entrants on all four CAID2
references. We are careful not to over-claim it: its margin over the CAID3
Disorder-PDB winner on the reported metric is $+0.0043$ at $p = 0.20$, and the
manuscript says so. What it does separate on is the calibration-invariant axis,
on all five benchmarks after Holm correction over thirty comparisons --- a
separation the benchmark's own metric provably cannot make.

The work has not been published elsewhere and is not under consideration by
another journal. All data are public, all code and the Lean development are
available, and the pre-registrations are committed with timestamps predating the
runs they govern. I have no competing interests.

Thank you for your consideration.

\closing{Yours sincerely,}

\end{letter}
\end{document}
